\documentclass[10pt,a4paper]{amsart}
\usepackage[margin=19mm]{geometry}
\usepackage{amsmath,amssymb,mathtools}
\usepackage[scr=rsfs,cal=euler]{mathalfa}
\usepackage{xfrac}
\usepackage[unicode,hidelinks,bookmarks=false]{hyperref}
\newcommand{\ie}{i.e.\ }
\newcommand{\cf}{cf.\ }
\newcommand{\resp}{respectively\ }
\newcommand{\Subsection}{\subsection}
\providecommand{\nobreakdash}{\nobreak\hbox{-}}
\setlength{\emergencystretch}{2em}
\multlinegap0pt
\usepackage{bm}
\usepackage{bbm}
\usepackage{dsfont}
\usepackage{xspace}
\usepackage{cancel}
\usepackage{mathtools}
\usepackage{amsthm}
\makeatletter
\@ifundefined{theorem}{\newtheorem{theorem}{Theorem}}{}
\@ifundefined{lem}{\newtheorem{lem}[theorem]{Lemma}}{}
\makeatother
\title[On the invariant surface area functionals in 3-dimensional C]{On the invariant surface area functionals in 3-dimensional CR geometry}
\author{Pak Tung Ho}
\date{Source: arXiv:2510.02632v2 (CC BY 4.0)}
\begin{document}
\maketitle

\noindent\emph{Source and licence.} On the invariant surface area functionals in 3-dimensional CR geometry. Version arXiv:2510.02632v2 (CC BY 4.0), distributed under the Creative Commons Attribution 4.0 International licence, as recorded in the arXiv licence metadata for this exact version and shown on the abstract page https://arxiv.org/abs/2510.02632v2. Full rights notice: licenses/arxiv-2510.02632-CC-BY-4.0.txt.\par\smallskip

\noindent\textit{Editorial scope.} Counted unit: the Lem whose source label is \texttt{lem5.2} in the article (source0.tex lines 1877--1883, source label \texttt{lem5.2}), delivered together with the article's own complete original proof of it (source0.tex lines 1884--1942, one \texttt{proof} environment of 59 lines). No mathematics is written, repaired or re-derived: statement and proof are the article's own text line for line. Required context retained uncounted: 5.1 (lines 1124--1129); 5.2 (lines 1134--1136); 5.9 (lines 1201--1215); 5.10 (lines 1220--1227); thm5.1 (lines 1230--1244); 10.3 (lines 1584--1589); 10.6 (lines 1621--1623); 10.6A (lines 1642--1644); 10.8 (lines 1654--1656); 10.17 (lines 1765--1775); 10.21 (lines 1807--1809). Every definition, display and label of those lines is retained unchanged. Omitted from this package and not redistributed: 1--1123, 1130--1133, 1137--1200, 1216--1219, 1228--1229, 1245--1583, 1590--1620, 1624--1641, 1645--1653, 1657--1764, 1776--1806, 1810--1876, 1943--3440. Nothing inside a retained excerpt is omitted or altered: every retained body is delivered exactly as the article's lines are written, so no omission receipt of any kind accompanies this package. Citations: the retained excerpts carry no citation command at all, so no bibliography entry of the article is transcribed and no reference is redistributed. Reference adaptation: none was needed and none was made; every reference inside the delivered unit resolves inside it, so no unresolved reference or citation is produced. Printed slips kept literal: no printed word, symbol, hypothesis or number of the counted statement or of its proof was altered. Layout and class: the article's own class is \texttt{amsart} with packages including graphicx, amsmath, amssymb, amsthm, esint, mathrsfs, amsfonts; the wrapper is the lane's pinned \texttt{amsart} editorial wrapper at 10pt on A4, which restates the theorem-like environments the retained text needs and adds the article's own macro definitions that the retained text needs (none), with extra wrapper packages (bm, bbm, dsfont, xspace, cancel, mathtools). The delivered numbering is the unit's own. Assets: no retained span contains a figure environment, a graphics-inclusion command, a PDF-page inclusion, a table or a TikZ picture; no image file is copied into this package, the package declares no asset dependency and no image file is redistributed with it. Licence: version arXiv:2510.02632v2 of this article is distributed under the Creative Commons Attribution 4.0 International licence, as recorded in the arXiv licence metadata for this exact version and shown on the abstract page https://arxiv.org/abs/2510.02632v2. A full rights notice is delivered at licenses/arxiv-2510.02632-CC-BY-4.0.txt. Characters: every retained excerpt is pure ASCII, so the delivered engine, XeLaTeX with its default Latin Modern text font, reports no missing character.
\begin{equation}\label{5.1}
\begin{split}
\mathcal{E}_1&=\frac{1}{2}e_1(|H_{cr}|^{1/2}\mathfrak{f})+\frac{3}{2}|H_{cr}|^{1/2}\alpha\mathfrak{f}\\
&\hspace{4mm}+\frac{1}{2}sign(H_{cr})|H_{cr}|^{1/2}\left(9h_{00}+6h_{11}h_{10}+\frac{2}{3}h_{11}^3\right).
\end{split}
\end{equation}

\begin{equation}\label{5.2}
H_{cr}=e_1(\alpha)+\frac{1}{2}\alpha^2-\mbox{Im} A_{11}+\frac{1}{4}W+\frac{1}{6}H^2. 
\end{equation}

\begin{equation}\label{5.9}
\begin{split}
|H_{cr}|\mathfrak{f}
&=\left(e_1(\alpha)+\frac{1}{2}\alpha^2+\frac{1}{3}H^2-\mbox{Im} A_{11}+\frac{1}{4}W\right)\Big(e_1(H)-2\alpha H+3 \mbox{Re} A^1_{\overline{1}}\Big)\\
&\hspace{4mm}
+H(T+\alpha e_2)(H)-3\alpha He_1(\alpha)-3\alpha^3H-\frac{3}{2}\alpha HW\\
&\hspace{4mm}+3\alpha H \mbox{Im} A^1_{\overline{1}}+3\alpha H \mbox{Im} A_{11}
+3H
\mbox{Re}\Big(\frac{1}{6}W^{,1}+\frac{2i}{3}(A^{11})_{,1}\Big)\\
&\hspace{4mm}
+\frac{3}{2}(T+\alpha e_2)\left(e_1(\alpha)+\frac{1}{2}\alpha^2-\mbox{Im} A_{11}+\frac{1}{4}W\right)+\frac{3}{2}\alpha H e_1(\alpha)+\frac{3}{2}\alpha^3 H\\
&\hspace{4mm}+\frac{3}{4}\alpha HW-\frac{3}{2}\alpha H \mbox{Im} A^1_{\overline{1}}-\frac{3}{2}\alpha H \mbox{Im}A_{11}
-\frac{3}{2}H\mbox{Re}\Big(\frac{1}{6}W^{,1}+\frac{2i}{3}(A^{11})_{,1}\Big).
\end{split}
\end{equation}

\begin{equation}\label{5.10}
\begin{split}
&9h_{00}+6h_{11}h_{10}+\frac{2}{3}h_{11}^3\\
&=9(T+\alpha e_2)(\alpha)
+9\mbox{Im}\Big(\frac{1}{6}W^{,1}+\frac{2i}{3}(A^{11})_{,1}\Big)-9\alpha \mbox{Re} A^1_{\overline{1}}\\
&\hspace{4mm}+6He_1(\alpha)+3H\alpha^2-6H\mbox{Im} A_{11}+\frac{3}{2}HW+\frac{2}{3}H^3.
\end{split}
\end{equation}

\begin{theorem}\label{thm5.1}
Let $\Sigma$ be a smooth surface in a $3$-dimensional 
pseudohermitian manifold $(M,J,\theta)$. 
Suppose $\Sigma$ is nonsingular and $H_{cr}\neq 0$ on $\Sigma$. 
Then $\Sigma$ satisfies the Euler-Lagrange equation 
for the energy functional $E_1$
if and only if $\mathcal{E}_1=0$, 
where $\mathcal{E}_1$ is given as in 
\eqref{5.1}, 
and 
$H_{cr}$, $\mathfrak{f}$
and $9h_{00}+6h_{11}h_{10}+\frac{2}{3}h_{11}^3$ 
are given as in \eqref{5.2}, \eqref{5.9}
\eqref{5.10} respectively. 
\end{theorem}

\begin{equation}\label{10.3}
\begin{split}
T&=i\left(z_1\frac{\partial}{\partial z_1}+z_2\frac{\partial}{\partial z_2}-\overline{z}_1\frac{\partial}{\partial\overline{z}_1}-\overline{z}_2\frac{\partial}{\partial\overline{z}_2}\right)\\
&=\frac{\partial}{\partial\varphi_1}+\frac{\partial}{\partial\varphi_2}.
\end{split}
\end{equation}

\begin{equation}\label{10.6}
A_{11}(t)=\overline{A_{\overline{1}\,\overline{1}}(t)}=\overline{\left(\frac{4it}{1-t^2}\right)}=-\frac{4it}{1-t^2},~~\omega(t)_1^1=-\frac{2i(1+t^2)}{1-t^2}\theta(t),~~W=\frac{2(1+t^2)}{1-t^2}.
\end{equation}

\begin{equation}\label{10.6A}
\alpha\equiv 0
\end{equation}

\begin{equation}\label{10.8}
e_1=c\left(-\frac{\rho_2}{\rho_1}\frac{\partial}{\partial\varphi_1}+\frac{\rho_1}{\rho_2}\frac{\partial}{\partial\varphi_2}\right)
\end{equation}

\begin{equation}\label{10.17}
\begin{split}
\theta\wedge e^1
&=|a|\rho_1\rho_2d\varphi_1\wedge d\varphi_2
+\frac{it(\overline{a}^2-a^2)}{(1-t^2)|a|}\Big[\rho_1^2(\rho_2d\varphi_1\wedge d\rho_1-\rho_1 d\varphi_1\wedge d\rho_2)\\
&\hspace{5cm}
+\rho_2^2(\rho_2 d\varphi_2\wedge d\rho_1-\rho_1d\varphi_2\wedge d\rho_2)\Big],\\
\theta\wedge e^1\wedge e^2
&=\rho_1 d\rho_1\wedge d\varphi_1\wedge d\varphi_2, 
\end{split}
\end{equation}

\begin{equation}\label{10.21}
H=-\frac{\rho_2^2-\rho_1^2}{\rho_2\rho_1}\left[|a|-\frac{t}{1-t^2}\left(4|a|+\frac{(\overline{a}^2-a^2)^2}{2|a|^3}\right)\right].
\end{equation}

\begin{lem}\label{lem5.2}
Suppose that $t\neq -4+\sqrt{15}$. 
Then the Clifford torus
 $\Sigma_{[\frac{\sqrt{2}}{2}]}$ 
in the Rossi sphere $S^3_{t}$
is a critical point for the energy $E_1$ with nonzero energy.
\end{lem}

\begin{proof}
It follows from  (\ref{5.2}), (\ref{5.9}),
(\ref{5.10}), (\ref{10.6}) and (\ref{10.6A})
that 
\begin{equation}\label{10.28}
\begin{split}
H_{cr}&=-\mbox{Im} A_{11}+\frac{1}{4}W+\frac{1}{6}H^2=\frac{1+8t+t^2}{2(1-t^2)}+\frac{1}{6}H^2,\\
|H_{cr}|\mathfrak{f}&=\left(\frac{1}{3}H^2-\mbox{Im} A_{11}+\frac{1}{4}W\right)e_1(H)+HT(H)\\
&=\left(\frac{1}{3}H^2+\frac{1+8t+t^2}{2(1-t^2)}\right)e_1(H)+HT(H),\\
9h_{00}+6h_{11}h_{10}+\frac{2}{3}h_{11}^3&=-6H\mbox{Im} A_{11}+\frac{3}{2}HW+\frac{2}{3}H^3
=\frac{6(1+8t+t^2)}{2(1-t^2)}H+\frac{2}{3}H^3.
\end{split}
\end{equation}
Note that it follows from (\ref{10.21}) that 
\begin{equation}\label{10.31}
H=0~~\mbox{ whenever }\rho_1=\rho_2=\frac{\sqrt{2}}{2}.
\end{equation}
On the other hand, from (\ref{10.3}) and (\ref{10.8}), 
we see that $e_1$ and $T$ only involve derivatives
with respect to $\varphi_1$ and $\varphi_2$. 
In particular, it follows from (\ref{10.21}) that 
\begin{equation*}
\begin{split}
e_1(H)&=-\frac{\rho_2^2-\rho_1^2}{\rho_2\rho_1}e_1\left[|a|-\frac{t}{1-t^2}\left(4|a|+\frac{(\overline{a}^2-a^2)^2}{2|a|^3}\right)\right],\\
T(H)&=-\frac{\rho_2^2-\rho_1^2}{\rho_2\rho_1}T\left[|a|-\frac{t}{1-t^2}\left(4|a|+\frac{(\overline{a}^2-a^2)^2}{2|a|^3}\right)\right],
\end{split}
\end{equation*}
which implies that 
\begin{equation}\label{10.30}
\begin{split}
e_1(H)=0~~\mbox{ and }~~
T(H)=0~~\mbox{ whenever }\rho_1=\rho_2=\frac{\sqrt{2}}{2}. 
\end{split}
\end{equation}
Similarly, we can see that 
\begin{equation}\label{10.30A}
X_1X_2\cdots X_n(H)=0~~\mbox{ whenever }\rho_1=\rho_2=\frac{\sqrt{2}}{2},
\end{equation}
as long as the vector fields $X_j=e_1$ or $T$ for all $1\leq j\leq n$. 
Combining 
(\ref{5.1}) and (\ref{10.31})-(\ref{10.30A}), we conclude that 
$\mathcal{E}_1=0$ 
whenever $\rho_1=\rho_2=\frac{\sqrt{2}}{2}$. 
Hence, it follows from Theorem \ref{thm5.1} that the Clifford torus
 $\Sigma_{[\frac{\sqrt{2}}{2}]}$ 
in the Rossi sphere $S^3_{t}$
is a critical point for the functional $E_1$. 
On the other hand, 
it follows from (\ref{10.17}), (\ref{10.28}) and (\ref{10.31})
that 
\begin{equation*}
\begin{split}
E(\Sigma_{[\frac{\sqrt{2}}{2}]})
&=\int_{\Sigma_{[\frac{\sqrt{2}}{2}]}}dA_1=\int_{\Sigma_{[\frac{\sqrt{2}}{2}]}}|H_{cr}|^{\frac{3}{2}}\theta\wedge e^1=\int_{\Sigma_{[\frac{\sqrt{2}}{2}]}}\left|\frac{1+8t+t^2}{2(1-t^2)}\right|^{\frac{3}{2}}\rho_1\rho_2
\end{split}
\end{equation*}
which is positive by the assumption on the values of $t$. 
This proves the assertion. 
\end{proof}

\end{document}
